\begin{table}[htbp]
\centering
\scriptsize
\caption{Test 2: client-input sensitivity for the illustrative retirement plan.}
\label{tab:gbi-client-sensitivity}
\begin{tabular}{llrrrrrr}
\toprule
Scenario & Risk & Assets & Contribution & Retirement gap & Projected wealth & Funding & Success \\
\midrule
Base client & Moderate & \$750,000 & \$25,000 & \$32,000 & \$1,216,109 & 105.7% & 66.6% \\
Conservative risk & Conservative & \$750,000 & \$25,000 & \$32,000 & \$1,038,574 & 90.3% & 8.5% \\
Assets -20\% & Moderate & \$600,000 & \$25,000 & \$32,000 & \$1,002,754 & 87.2% & 11.6% \\
Contribution \$15k & Moderate & \$750,000 & \$15,000 & \$32,000 & \$1,156,375 & 100.6% & 48.3% \\
Essential expenses +\$15k & Moderate & \$750,000 & \$25,000 & \$47,000 & \$1,216,109 & 105.7% & 66.6% \\
\bottomrule
\end{tabular}
\vspace{2pt}
\begin{minipage}{0.97\linewidth}
\scriptsize
Notes: All scenarios retain a five-year horizon and a \$1.15 million target. The expense scenario changes the staged retirement cash-flow gap; the current terminal-wealth bootstrap does not yet deduct retirement spending, so its success probability remains unchanged.
\end{minipage}
\end{table}
